%Komplexe Wechselstromrechnung

\newvideofile{Wechsel}{Wechselstromrechnung}

\begin{frame}
    \fta{Complex alternating current calculation}
    
    \s{
        Direct current networks use constant voltage sources and current sources. In contrast, 
        complex alternating current calculations use sources with sinusoidal alternating quantities. The electrical analysis of components 
        in an alternating voltage network is complex.
        The complex voltages \underline{U} and currents \underline{I} generate time-dependent voltage values u(t) and current values 
        i(t). Equation \ref{GleichungURIKomplex} describes the specified relationship between complex and time-dependent
        current and voltage values for an ohmic resistor.
        
        Complex and time-dependent current and voltage specifications in AC calculation:
    
    
    \begin{eq}
        \underline{U} = \underline{Z}_\mathrm{R} \cdot \underline{I} \quad \rightarrow \quad u(t) = R \cdot i(t)  \label{GleichungURIKomplex}
    \end{eq}
    
    
        Apart from the complexity of alternating current calculations, the rules for calculating electrical networks still apply,
        such as mesh and node rules, the behaviour of series and parallel circuits, node potential analysis, and 
        mesh current analysis. In all these examples, care must be taken to calculate with complex alternating quantities. For the 
        explanation of electrical networks with complex alternating quantities, the following topics are examined in more detail below:
        For the calculation of electrical networks with alternating quantities, the following topics are examined in more detail below:
    }
    
    
    \begin{Learning objectives}{Alternating current calculation}
        The students
        \begin{itemize}
            \item know the complex impedance and complex admittance.
            \item understand the behaviour of a resistor, a capacitor and an inductor at an alternating voltage.
            \item know alternating current sources and the associated conversion rules.
        \end{itemize}
    \end{Learning objectives}
    
    \speech{WechselLrnz}{1}{Wie auch bei Gleichstromnetzwerken wird bei der Wechselstromrechnung mit einem elektrichen Widerstand gearbeitet.
    Jedoch ist der elektrische Widerstand wie schon die Größen von Strom und Spannung bei der Wechselstrombetrachtung eine komplexe Größe.
    Hier soll nun folgend erläutert werden, was die komplexe Impedanz und die kompolexe Attmitanz ist.
    Das Verständnis über das Verhalten von Bauelementen an einer Wechselspannung vermittelt werden 
    und gezeigt werden, wie Wechselquellen ineinander umgewandeln werden können}
\end{frame}

\begin{frame}
    \renewcommand{\figurename}{Figure}
    \ftb{Impedance and admittance in the complex plane}
    
    \s{
        A phase shift between the voltage and the current results in the generation of a reactive component of the impedance. However, the reactive component cannot be represented 
        on the real axis, which is why the real axis is extended by an imaginary axis, which together form the 
        complex plane. Figure \ref{BildKomplexeImpedanz} shows the phasor of the complex impedance in the complex plane 
        with imaginary and real components. 
        
        \fu{
            \input{Tikz/src/Komplexe_Wechselstromrechnung/Komplexe_Impedanz}
        }{{\bf Complex impedance pointer.} The impedance in the complex plane with the division into real part R and imaginary part X. \label{BildKomplexeImpedanz}}
        
    }
    
    \b{
        \begin{minipage}[t]{0.4\textwidth}
            \begin{itemize}
                \item<1-> Complex pointer of impedance
                \item<2-> Active component of impedance
                \item<3-> Reactive component of impedance
                \item<4-> Active resistance and reactive resistance of impedance
            \end{itemize}
        \end{minipage}
        \begin{minipage}[t]{0.4\textwidth}
            \fu{
                \begin{tikzpicture}
                    \draw (0,0) coordinate (U);
                    \draw[very thin,color=gray] (-0.1,-0.1) grid (3.2,3.1);
                    \draw[->] (-0.2,0) -- (3.4,0) node[right] {$\Re(\underline{Z})$};
                    \draw[->] (0,-0.1) -- (0,3.2) node[above] {$\Im(\underline{Z})$};
                    \draw (3,0) coordinate (R);
                    \draw (3,3) coordinate (Z);
                    \draw[->] (0,0) -- (2.5,2.5);
                    \draw(1.25,1.75) node {$\hat{Z}$};
                    \draw pic[draw, angle radius = 1cm, ->] {angle = R--U--Z};
                    \draw(1.2,0.5) node {$\varphi_Z$};
                    \draw(3,2.5) node [right] {$\underline{Z}=\hat{Z} \cdot e^{j \varphi_\mathrm{Z}}$};
                    \pause
                    
                    \draw[dashed](2.5,2.5) -- (2.5,-0.2) node[below] {$R$};
                    \pause
                    
                    \draw[dashed](2.5,2.5) -- (-0.2,2.5) node[left] {$j \cdot X$};
                \end{tikzpicture}
            }{}
        \end{minipage}
    }
    
    \s{
        According to equation \ref{GleichungImpedanz}, the complex impedance \underline{Z} also consists of a real part and an imaginary part. The
        real part consists of the ohmic component, which represents the effective resistance of an impedance. The effective resistance is also referred to as 
        resistance R. The second component of impedance is the imaginary part, which is also referred to as reactance X.
        The reactance of an electrical network, for example, consists of the capacitive and inductive effects
        and is assigned the imaginary unit j. As in direct current technology, the unit of complex impedance is the ohm $\Omega$. 
    }
    
    \only<4>{
        \begin{eq}
            Impedance = Resistance + j \cdot Reactance \quad \rightarrow \quad \underline{Z} = R + j \cdot X   \label{GleichungImpedanz}
        \end{eq}
        \begin{eq}
            [\underline{Z}] = 1\ Ohm = 1\ \Omega \nonumber
        \end{eq}
    }

    \speech{WechselImp}{1}{Hier wird der komplexe Zeiger der Impedanz Zett gezeigt.
    Die Impedanz lässt sich ebenso wie die anderen komplexen Größen in einen Realteil und einen Imaginärteil unterteilen.}
    \speech{WechselImp}{2}{Bei der Impedanz wird der Realteil auch als Wirkanteil R bezeichnet.}
    \speech{WechselImp}{3}{Der Imaginärteil der Impedanz als Blindanteil X.}
    \speech{WechselImp}{4}{Die komplexe Impedanz setzt sich so aus dem Wirkanteil, der Resistanz, 
    und dem Blindanteil, der Reaktanz, zusammen.
    Die Einheit der komplexen Impedanz Zett ist das Ohm.}
\end{frame}




\begin{frame}
    \ftx{Impedance and admittance in the complex plane}
    
    \s{
        As in direct current technology, there is a conductance value for electrical resistance in complex alternating current technology.
        This reciprocal value of the complex impedance is referred to as admittance \underline{Y} (see equation \ref{GleichungImpedanzAttmitanz}). The unit of the complex conductance remains 
        Siemens S, as in direct current technology.
    }
    
    \b{
        Complex conductance as the reciprocal of complex impedance:
    }
    
    \begin{eq}
        R = \frac{1}{G} \quad \rightarrow \quad \underline{Z} = \frac{1}{\underline{Y}}   \label{GleichungImpedanzAttmitanz}
    \end{eq}
    
    \s{ 
        Admittance can also be differentiated into a real conductance and a reactive conductance. 
        The real conductance G of the admittance is referred to as conductance, and the reactive conductance B of the 
        admittance is referred to as susceptance (see Equation \ref{GleichungAdmittanz} ).
    }
    
    \onslide<2->{
    \b{
        Effective conductance and reactive conductance of admittance:
    }
    
    \begin{eq}
        Admittance = Conductance + j \cdot Susceptibility \quad \rightarrow \quad \underline{Y} = G + j \cdot B           \label{GleichungAdmittanz}
    \end{eq}
    
    \begin{eq}
        [\underline{Y}] = 1\ Siemens = 1\ S \nonumber
    \end{eq}}
    
    \speech{WechselAtt}{1}{Bei der Betrachtung von Gleichstromnetzwerken,
    ergibt sich der Leitwert aus dem Kehrwert des Widerstandes.
    Der komplexe Leitwert, die Attmitanz Y, 
    ergibt sich aus dem Kehrwert der komplexen Impedanz.}
    \speech{WechselAtt}{2}{Auch hier wird die Unterteilung in Wirkleitwert und Blindleitwert vorgenommen.
    Die Attmitanz setzt sich so aus der Konduktanz G und der Suseptanz B zusammen.
    Die Einheit der komplexen Admittanz Y ist das Siemens.}
\end{frame}


\begin{frame}
    \ftx{Complex impedance and complex admittance}
    
    \begin{Key point}{Complex impedance and complex admittance}      
        The direct current resistance becomes the complex impedance \underline{Z} for an alternating voltage. The 
        reciprocal of the resistance, the conductance, becomes the complex admittance \underline{Y} for an alternating voltage.    
    \end{Key point}

    \speech{WechselMrk1}{1}{Der Gleichstromwiderstand wird bei einer Wechselspannun zur komplexen Impedanz Z. 
    Der Kehrwert des Widerstandes, der Leitwert, wird bei einer Wechselspannung zur komplexen Admittanz Y.
    Zur Erinnerung.
    Komplexe Größen werden durch einen Unterstrich gekennzeichnet.}
\end{frame}


\begin{frame}
    \ftb{Complex resistance}
    
    \s{
        The complex resistance of an ohmic consumer consists of the active component R and the reactive component of the ohmic resistance.
        The ohmic resistance R embodies the property of conductive materials to convert electrical power into thermal power.
        The reactive component of the resistance does not convert the power of alternating currents into heat.         
        The reactive component of the electrical resistance is denoted as $X_\mathrm{R}$. However, since the ideal electrical resistance is purely
        resistive and therefore purely real, the symbol R is often used. Equation \ref{GleichungWid1}
        explains these relationships.
    }
    
    \b{
        \begin{itemize}
            \item<1-> Complex resistance at an ohmic load
            \item<2-> No reactive component in an ideal ohmic resistance
            \item<3-> Ideal complex resistance as pure active component
        \end{itemize}
    }
    
    \begin{eq}
        \underline{Z}_\mathrm{R} = R + j \cdot X_\mathrm{R} \onslide<2->{\quad \rightarrow \quad X_\mathrm{R-ideal} = j \cdot 0\ \Omega} \onslide<3->{\quad \rightarrow \quad \underline{Z}_\mathrm{R} = R}       \label{GleichungWid1}
    \end{eq}
    
    \s{
        The time-dependent values for voltage and current result in the relationships of ideal ohmic resistance 
        according to equation \ref{GleichungWid2}.
    }
    
    \onslide<4->{
    \begin{eq}
        \underline{U} = \underline{Z}_\mathrm{R} \cdot \underline{I} \quad \rightarrow \quad u(t) = R \cdot i(t)        \label{GleichungWid2}
    \end{eq}}
    
    \speech{WechselWiderErkl}{1}{Die Impedanz an einem rein ohmschen Verbraucher wird durch Zett R beschrieben.}
    \speech{WechselWiderErkl}{2}{Ist ein Verbraucher rein Ohmsch, bewirkt er keinen Blindanteil. 
    X R wird somit zu Null und das Produkt der Reaktanz entfällt.}
    \speech{WechselWiderErkl}{3}{Der ideale Widerstand erfolgt auf diese Weise als reiner Wirkwiderstand}
    \speech{WechselWiderErkl}{4}{Das ohmsche Gesetz }
\end{frame}

\begin{frame}
    \renewcommand{\figurename}{Figure}
    \ftx{Complex resistance}
    
    \s{
        An ideal electrical resistance is a pure active resistance and has no reactive resistance components.
        Real electrical resistance is subject to parasitic effects. The circuit diagrams of an ideal and a 
        real electrical resistance are compared in Figure \ref{BildRealerWiderstand}. These parasitic 
        effects have an inductive effect, illustrated by the coil, and a capacitive effect, illustrated by the capacitance, on the 
        circuit. In addition to the frequency dependence of the parasitic effects, the voltage and current are also divided in such a way
        that the entire power no longer drops across the resistor. The impedance of the ideal electrical resistor is not 
        frequency-dependent.
    }
    
    \b{
        Electrical resistance in alternating current calculations
        \begin{itemize}
            \item<1-> Ideal resistance
            \item<2-> Real resistance
            \item<3-> No phase shift with ideal resistance
        \end{itemize}
    }
    
    \fu{
        \input{Tikz/src/Komplexe_Wechselstromrechnung/Gegenueberstellung_idealer_realer_Widerstand}
    }{{\bf Comparison of ideal and real electrical resistance.} Unlike ideal resistance, real resistance exhibits
        serial inductive and parallel capacitive effects. \label{BildRealerWiderstand}}

    \speech{WechselWiderBauteil}{1}{In der Realität lässt sich diese vereinfachte Darstellung eines ohmschen Widerstandes leider nicht in gänze Nutzen.}
    \speech{WechselWiderBauteil}{2}{Beispielsweise sorgen die Zu und Ableitungen an den Bauteilen in der Pärepherie zu parasitären Effekten, 
    welche durch das Schaltbild eines Realen Widerstandes abgebildet werden.
    Die parasitären Effekte lassen sich durch die Induktivitäten von Spulen und Kapazitäten von Kondensatoren erklären.}
    \speech{WechselWiderBauteil}{3}{Wird der ohmsche Widerstand aber als rein ideales Bauteil angesehen, erfolgt keine Phasenverschiebung zwischen Strom und Spannung.}
\end{frame}


\begin{frame}
    \renewcommand{\figurename}{Figure}
    \ftx{Complex resistance}
    
    \s{
        With ideal electrical resistance, there is no phase shift between the voltage and the current. In 
        Figure \ref{BildPhasenWiderstand}, the sine waves for the voltage and current are plotted on an ideal resistor. 
        They lie on top of each other without any phase shift, meaning they are in phase.
    }
    
    \b{
        Phase shift at ideal ohmic resistance:
        \begin{itemize}
            \item<1-> Voltage $u_\mathrm{R}(t)$
            \item<2-> Current $i_\mathrm{R}(t)$
            \item<3-> Voltage and current at the ideal ohmic resistance without phase shift
        \end{itemize}
    }
    
    \fu{
        \input{Tikz/src/Komplexe_Wechselstromrechnung/Sinunsschwingung_Widerstand}
    }{{\bf Sine waves on a resistor.} Alternating voltage (blue) and alternating current (red) on an ideal resistor. This results in
        no phase shift between the alternating voltage and the alternating current at the resistor. \label{BildPhasenWiderstand}}

    \speech{WechselWiderPhase}{1}{Die Phasenverschiebung soll trotzdem noch einmal näher betrachtet werden, da sie bei den anderen Bauteilen noch wichtig wird.
    Dargestellt wird die zeitabhängige Spannung U von T.}
    \speech{WechselWiderPhase}{2}{Eben auf der Spannung liegt der zeitliche Verlauf des Stromes I von T.
    Die Spannung und der Strom unterscheiden sich hier in ihrer Amplitude, jedoch nicht in ihrer Phase.}
    \speech{WechselWiderPhase}{3}{Die Grafische Darstellung zeigt so, 
    dass die Differenz zwischen den Phasenwinkeln von Spannung und Strom null ergibt.}
\end{frame}


\begin{frame}
    \ftx{Sinusoidal oscillation at a resistor}
    
    \begin{Key point}{Sinusoidal oscillation at a resistor}
        The {\bf electrical resistance} indicates between the alternating voltage and the corresponding alternating current {\bf no
        phase shift}.     
    \end{Key point}

    \speech{WechselMrk2Wid}{1}{Der elektrische Widerstand weist zwischen der Wechselspannung 
    und dem Wechselstrom keine Phasenverschiebung auf.}
\end{frame}


\begin{frame}
    \ftb{Capacity}
    
    \s{
        As a component in electrical engineering, the capacitor has the ability to store energy between the two plates in 
        the electric field. The size of a capacitor is specified in terms of its capacitance C. The unit 
        of capacitance is the farad F. The complex impedance of a capacitor can be calculated using equation \ref{GleichungKap1}. 
        In contrast to the ideal resistor, the ideal capacitor in an alternating current circuit does not act solely as an
        effective resistance. The real capacitor acts both through an effective component and through an reactive component. The 
        circular frequency $\omega$ is frequency-dependent. This results in low impedance values for the 
        capacitor at high frequencies. Conversely, high impedance values result at low frequencies. The reactive component of the capacitor
        acts as capacitive reactance.
        The capacitor can also be calculated using complex admittance. The complex admittance is derived from
        the reciprocal of the complex impedance of the capacitor. 
    }
    
    \b{
        Negative reactive resistance $\leftrightarrow$ Representation in the negative imaginary domain
    }
    
    \begin{eq}
        \underline{Z}_\mathrm{C} = \frac{1}{j \omega C} \quad \rightarrow \quad \underline{Y}_\mathrm{C} = j \omega C          \label{GleichungKap1}
    \end{eq}
    
    \s{
        The ideal capacitor acts purely in the reactive component, so according to equation \ref{GleichungKap2} from the complex notation $Z_C$, the 
        reactive component of the capacitor $X_C$ with time-dependent voltage and current.
    }
    
    \onslide<2->{
    \b{
        Ideal capacitor for alternating voltage and alternating current:  
    }
    
    \begin{eq}
        \underline{U} = \underline{Z}_\mathrm{C} \cdot \underline{I} \qquad \qquad u(t) = \underline{X}_\mathrm{C} \cdot i(t)     \label{GleichungKap2}
    \end{eq}}
    
    \s{
        It is assumed that the alternating voltage at a capacitor has a phase shift of 0°. 
        With this assumption, the phase shift of the current at the capacitor can be determined using equation \ref{GleichungKap3}. 
        Due to the phase shift of the current, the curve of a cosine function contrasts with the sine function
        of the voltage. If both the voltage and the current are to be described by the sine function, the 
        phase shift $\Delta\varphi$ of $\pi/2$ is included in the argument. 
    }
    
    \onslide<3->{
    \b{
        Phase shift of the current at the ideal capacitor: 
    }
    
    \begin{eq}
        u(t) = \hat{U} \cdot \sin(\omega t)      \qquad      i(t) = \hat{I} \cdot \cos(\omega t) = \hat{I} \cdot \sin(\omega t + \pi/2)        \label{GleichungKap3}
    \end{eq}}
    
    \speech{WechselKapaErkl}{1}{Die Impedanz einer Kapazität wird aus dem Ausdruck von Zett C gebildet.
    Hierbei steht der Ausdruck jot Omega C im Nenner.
    Zur Berechnung der kapazitiven Attmitanz wird der Kehrwert aus der kapazitiven Impedanz berechnet.
    Ein kapazitiver Blindwiderstand wird im negativen Imaginärbereich dargestellt.}
    \speech{WechselKapaErkl}{2}{Ein idealer kondensator besteht rein aus den kapazitiven Blindanteil der Impedanz.}
    \speech{WechselKapaErkl}{3}{An einer Kapazität kommt es zu einer Phasenverschiebung zwischen dem Strom und der Spannung.
    }
\end{frame}

\begin{frame}
    \renewcommand{\figurename}{Figure}
    \ftx{Capacity}
    
    \s{
        With inductors and capacitors, there is a phase shift between the alternating voltage and the
        alternating current. With capacitors, this means that the current leads the voltage by 90° (degrees). The described 
        phenomenon of phase shift is illustrated in Figure \ref{BildPhasenKondensator}. The current shown in red 
        at a capacitor reaches its maximum value $2\pi$ (radians) earlier than the blue-coloured 
        maximum value of the voltage.  
    }
    
    \b{
        Phase shift at the capacitor:
        \begin{itemize}
            \item Voltage $u_\mathrm{C}(t)$
            \item<2-> Current $i_\mathrm{C}(t)$
            \item<3-> Phase shift between current and voltage
            \item<4-> Phase shift $\Delta\varphi$
        \end{itemize}
    }
    
    \fu{
        \input{Tikz/src/Komplexe_Wechselstromrechnung/Phasenverschiebung_Kondensator}
    }{{\bf Phase shift between current and voltage at a capacitor.} The current leads the voltage across the capacitor by 90°.
        \label{BildPhasenKondensator}}
    
    \speech{WechselKapaPhase}{1}{Hier wird nun der Verlauf der Wechselspannung U C an einer Kapazität dargestellt. 
    Die Wechselspannung startet im Ursprung des Grafen und hat eine Halbwelle nach Pi absolviert.}
    \speech{WechselKapaPhase}{2}{Neben der Wechselspannung wird nun in Rot der Wechselstrom an der Kapazität dargstellt. 
    Er startet nicht im Ursprung. Die positive Halbwelle startet um Pi halbe Phsenverschoben im Vergleich zur Wechselspannun.}
    \speech{WechselKapaPhase}{3}{Der Verlauf des Wechselstromes erreicht den Maximalpunkt bei zwei Pi früher als der Verlauf der Wechselspannung.}
    \speech{WechselKapaPhase}{4}{Diese Phasendifferenz zwischen den beiden Phasenlagen beträgt bei einer Kapazität Pi halbe oder im Gradmaß 90 Grad.}
\end{frame}


\begin{frame}
    \ftx{Sinusoidal oscillation on a capacitor}
    
    \begin{Key point}{Sinusoidal oscillation on a capacitor}
        At the ideal capacity, the alternating current of the alternating voltage leads by 90°.     
    \end{Key point}

    \speech{WechselMrk3Kapa}{1}{An der idealen Kapazität eilt der Wechselstrom der Wechselspannung um 90 Grad vor.}
\end{frame}


\begin{frame}
    \renewcommand{\figurename}{Figure}
    \ftx{Ideal and real capacitor}
    
    \s{
        As a rule, calculations are based on ideal components and their properties. However, real capacitors, like 
        the real resistors described above, exhibit inductive and ohmic parasitic effects. The ideal capacitor is compared with the real capacitor in 
        Figure \ref{BildRealerKondensator}. The input and output lines exhibit inductive 
        ($L_\mathrm{paras.}$) and ohmic ($R_\mathrm{paras.2}$) components. In addition, the capacitor is subject to constant 
        self-discharge. This self-discharge occurs symbolically via the parallel resistance $R_\mathrm{paras.1}$.
    }
    
    \b{
        Ideal and real capacitor:
        \begin{itemize}
            \item Ideal capacitor
            \item<2-> Real capacitor with parasitic effects
        \end{itemize}
    }
    
    \fu{
        \input{Tikz/src/Komplexe_Wechselstromrechnung/Gegenueberstellung_idealer_realer_Kondensator}
    }{{\bf Comparison of an ideal and a real capacitor.} Unlike an ideal resistor, a real capacitor exhibits
        serial inductive and parallel capacitive effects. \label{BildRealerKondensator}}
    
    \speech{WechselKapaBauteil}{1}{}
    \speech{WechselKapaBauteil}{2}{}
\end{frame}


\begin{frame}
    \ftb{Inductance}
    
    \s{
        The coil is also an energy storage device, just like the capacitor. However, in the case of the coil, 
        the energy is not stored in an electric field, but in a magnetic field according to the laws of induction. The size of a coil 
        is given by its inductance L. The unit of inductance is Henry H. The complex impedance of a coil
        is given by the imaginary unit j, the angular frequency $\omega$ and the inductance L (see Equation \ref{GleichungInd1}).
        The complex admittance is calculated by the reciprocal of the coil impedance. This leads to relationships that develop contrary to
        the impedance calculations for capacitances. At high frequencies, the impedances of coils also become large in accordance with their 
        inductance and correspondingly small at low frequencies. 
    }
    
    \b{
        Positive reactance $\leftrightarrow$ Representation in the positive imaginary domain
    }
    
    \begin{eq}
        \underline{Z}_\mathrm{L} = j \omega L \quad \rightarrow \quad \underline{Y}_\mathrm{L} = \frac{1}{j \omega L}     \label{GleichungInd1}
    \end{eq}
    
    \s{
        As with the capacitor, the ideal coil acts purely as a reactive component of the impedance. In equation \ref{GleichungInd2},
        the complex notation of the coil impedance $Z_L$ is used on the left-hand side. On the right-hand side, the time function
        is represented solely by the notation of the reactive component of the coil $X_L$.
    }
    
    \b{
        Ideal coil for alternating voltage and alternating current:  
    }
    
    \begin{eq}
        \underline{U} = \underline{Z}_\mathrm{L} \cdot \underline{I} \qquad  \qquad u(t) = \underline{X}_\mathrm{L} \cdot i(t)        \label{GleichungInd2}
    \end{eq}
    
    \s{
        In principle, the time functions of inductances for voltage and current behave similarly to the time functions
        of capacitances. Only the phase shift of the current in coils is opposite to the phase shift of capacitors. 
        The phase shift $\Delta\varphi$ between the voltage and the current at an inductance is $\pi/2$.
    }
    
    \b{
        Phase shift of the current at the ideal coil:
    }
    
    \begin{eq}
        u(t) = \hat{U} \cdot \cos(\omega t) = \hat{U} \cdot \sin(\omega t + \pi/2)    \qquad     i(t) = \hat{I} \cdot \sin(\omega t)      \label{GleichungInd3}
    \end{eq}
    
\end{frame}

\begin{frame}
    \renewcommand{\figurename}{Figure}
    \ftx{Inductance}
    
    \s{
        Equivalent to the phase shift at a capacitor, there is also a phase shift at the inductance.
        However, this phase shift of the current at the inductance is opposite to the phase shift of the capacitance.
        In the specific case shown in Figure \ref{BildPhasenSpule}, it can be seen that the current (red) at an inductance lags behind the 
        voltage (blue).
    }
    
    \b{
        Phase shift at the coil:
        \begin{itemize}
            \item Voltage $u_\mathrm{L}(t)$
            \item<2-> Current $i_\mathrm{L}(t)$
            \item<3-> Phase shift between current and voltage
            \item<4-> Phase shift $\Delta\varphi$
        \end{itemize}
    }
    
    \fu{
        \input{Tikz/src/Komplexe_Wechselstromrechnung/Phasenverschiebung_Spule}
    }{{\bf Phase shift between current and voltage at a coil.} The current lags the voltage at the coil by 90°.
        \label{BildPhasenSpule}}
    
\end{frame}



\begin{frame}
    \ftx{Sinusoidal oscillation on a coil}
    
    \begin{Key point}{Sinusoidal oscillation on a coil}
        With ideal inductance, the alternating current lags the alternating voltage by 90°.     
    \end{Key point}
\end{frame}


\begin{frame}
    \renewcommand{\figurename}{Figure}
    \ftx{Ideale und reale Spule}
    
    \s{
        Like ohmic resistance and capacitors, real coils exhibit parasitic effects. The coil wire 
        has an ohmic component ($R_\mathrm{paras.}$) that affects the complex impedance of the coil. In addition, the winding 
        of the coil generates capacitive effects ($C_\mathrm{paras.}$). Figure \ref{BildRealeSpule} shows the equivalent circuit diagram of an 
        ideal coil and a real coil.
    }
    
    \b{
        Ideal and real coil:
        \begin{itemize}
            \item Ideal coil
            \item<2-> Real coil with parasitic effects
        \end{itemize}
    }
    
    \fu{
        \input{Tikz/src/Komplexe_Wechselstromrechnung/Gegenueberstellung_ideale_reale_Spule}
    }{{\bf Comparison of an ideal and a real coil.} The real coil exhibits serial ohmic and parallel
        capacitive parasitic effects. \label{BildRealeSpule}}
    
\end{frame}


\begin{frame}
    \renewcommand{\figurename}{Figure}
    \ftb{Sources of alternating quantities}
    
    \s{
        As with sources of direct current and direct voltage, there are alternating current sources with comparable characteristics. In an
        ideal alternating voltage source, the voltage is fixed. The alternating voltage does not change depending on 
        the current flowing through it. 
        The ideal alternating current source provides a fixed current . 
        This is not influenced by the connected load. The circuit symbols for an ideal 
        alternating voltage source and an ideal alternating current source are shown in Figure \ref{BildIdealeQuellen}. 
    }
    
    \b{
        Sources of alternating quantities:
        \begin{itemize}
            \item Ideal alternating voltage source
            \item<2-> Ideal alternating current source
        \end{itemize}
    }
    
    \fu{
        \input{Tikz/src/Komplexe_Wechselstromrechnung/Quellen_Wechselgroessen}
    }{{\bf Sources of alternating quantities.} Ideal AC voltage source and ideal AC current source. \label{BildIdealeQuellen}}
\end{frame}


\begin{frame}
    \renewcommand{\figurename}{Figure}
    \ftx{Real AC voltage source}
    
    \s{
        The two alternating current sources presented here are ideal sources. Real alternating current sources, like 
        sources of direct voltage and direct current, have internal resistance. This internal resistance manifests itself in sources of 
        alternating quantities as complex impedances. The real alternating voltage source therefore has a complex internal impedance
        $\underline{Z}_\mathrm{i}$, which is connected in series with an ideal alternating voltage source. In real 
        alternating voltage sources, the output voltage $\underline{U}$ depends on the connected load. The comparison
        between an ideal and a real alternating voltage source is shown in Figure \ref{BildIdealSpannung}.
    }
    
    \b{
        Alternating voltage source:
        \begin{itemize}
            \item Ideal AC voltage source
            \item<2-> Real AC voltage source with internal impedance $\underline{Z}_\mathrm{i}$
        \end{itemize}
    }
    
    \fu{
        \input{Tikz/src/Komplexe_Wechselstromrechnung/Gegenueberstellung_ideale_reale_Wechselspannungsquelle}
    }{{\bf Comparison of an ideal and a real alternating voltage source.} The real alternating voltage source has a
        series impedance $\underline{Z}_\mathrm{i}$. \label{BildIdealSpannung}}
    
    \s{
        The output voltage of a real AC voltage source can be determined using equation \ref{GleichungQuel1}.
        Here, the AC voltage of the AC voltage source is not directly applied to the terminals. It is reduced by the voltage
        across the internal impedance $Z_\mathrm{i}$. 
    }
    
    \only<3>{
        \b{
            Output voltage of a real AC voltage source:
        }
        
        \begin{eq}
            \underline{U} = \underline{U}_\mathrm{0} - \underline{Z}_\mathrm{i} \cdot \underline{I}     \label{GleichungQuel1}
        \end{eq}
    }
\end{frame}

\begin{frame}
    \ftx{Characteristics of real AC voltage sources}
    
    \s{
        Real AC voltage sources exhibit certain characteristics when unloaded, at no load and in the event of a short circuit. 
        When idle, there is an infinite impedance between the terminals. No closed circuit is formed 
        and no current flows (equation \ref{GleichungQuel2}). The entire voltage of the AC source 
        is applied to the open terminals, see equation \ref{GleichungQuel3}.
    }
    
    \b{
        No-load:
    }
    
    \begin{eq}
        \underline{I} = 0 \qquad \rightarrow \qquad \underline{U}_\mathrm{i} = \underline{I} \cdot \underline{Z}_\mathrm{i} = 0 \label{GleichungQuel2}
    \end{eq}
    \begin{eq}
        \underline{U}_\mathrm{0} - \underline{U}_\mathrm{i} - \underline{U} = 0 \qquad \rightarrow \qquad \underline{U} = \underline{U}_\mathrm{0} \label{GleichungQuel3}
    \end{eq}
    
    \s{
        In the event of a short circuit, the impedance between the terminals approaches zero and there is no voltage drop. This results in the 
        maximum possible short-circuit current, which is limited only by the internal impedance (equation \ref{GleichungQuel4}).
        The voltage of the AC voltage source thus acts in its entirety across the internal impedance (Equation \ref{GleichungQuel5}). The 
        short-circuit current can be determined according to Equation \ref{GleichungQuel6} from the ratio of the voltage across the internal impedance 
        $\underline{U}_\mathrm{i}$ and the internal impedance $\underline{Z}_\mathrm{i}$.
    }
    
    \b{
        Short circuit:
    }
    
    \begin{eq}
        \underline{U} = 0 \qquad \rightarrow \qquad \underline{I} = \underline{I}_\mathrm{k} \label{GleichungQuel4}
    \end{eq}
    \begin{eq}
        \underline{U}_\mathrm{0} - \underline{U}_\mathrm{i} - \underline{U} = 0 \qquad \rightarrow \qquad \underline{U}_\mathrm{i} = \underline{U}_\mathrm{0} \label{GleichungQuel5}
    \end{eq}
    \begin{eq}
        \underline{I}_\mathrm{k} = \frac{\underline{U}_\mathrm{i}}{\underline{Z}_\mathrm{i}} \label{GleichungQuel6}
    \end{eq}
    
\end{frame}


\begin{frame}
    \renewcommand{\figurename}{Figure}
    \ftx{Example task}
    \begin{expl}{Real Voltage source}{Real voltage source:}
        
        \begin{enumerate}
            \only<1>{\item[]	
                    A real voltage source (230 V, 50 Hz) has an internal impedance $\underline{Z}_\mathrm{i}$ as shown in the figure below, which consists of 
                    a coil $L=20\ mH$ and an ohmic resistance $R=10\ m\Omega$.
                   
                  \fu{
                      \resizebox{0.5\textwidth}{!}{
                          \input{Tikz/src/Komplexe_Wechselstromrechnung/Beispiel_reale_Spannungsquelle}
                      }
                  }{{\bf Example.} Alternating voltage source with an internal impedance consisting of a coil and a resistor. \label{BildBeispielQuelle}}
                  
                  Determine the {\bf open-circuit voltage} and the {\bf short-circuit current} AC voltage source.
                  }
                  
                  \only<2>{
            \item[] 
                  \onslide<2>{
                  \begin{eqa}
                      \underline{U}_\mathrm{0} &= 230\ V \cdot e^{\mathrm{j}0^o}          \nonumber    \\
                      \underline{Z}_\mathrm{i} &= R + j X_L       \nonumber    \\
                      \underline{Z}_\mathrm{i} &= 0,01\ \Omega + j \cdot 2\pi \cdot 50\ Hz \cdot 20\ mH      \nonumber    \\
                      \underline{Z}_\mathrm{i} &= 0,01\ \Omega + j 6,28\ \Omega       \nonumber
                  \end{eqa}
                  
                  {\bf Open-circuit voltage}:
                  \begin{eqa}
                      \underline{U} &= \underline{U}_\mathrm{0}        \nonumber    \\
                      \underline{U} &= 230\ V \cdot e^{\mathrm{j}0^o}         \nonumber
                  \end{eqa}
                  
                  {\bf Short-circuit current}:
                  \begin{eqa}
                      \underline{I}_\mathrm{k} &= \frac{\underline{U}_\mathrm{0}}{\underline{Z}_\mathrm{i}}           \nonumber        \\
                      \underline{I}_\mathrm{k} &= \frac{230\ V \cdot e^{\mathrm{j}0^o}}{0,01\ \Omega + j 6,28\ \Omega}           \nonumber        \\
                      \underline{I}_\mathrm{k} &= (0,058 - j36,6)\ A = 36,6\ A \cdot e^{-\mathrm{j}89,9^o}    \nonumber
                  \end{eqa}
                  }
                  }    
        \end{enumerate}
    \end{expl}
\end{frame}


\begin{frame}
    \renewcommand{\figurename}{Figure}
    \ftx{Real AC source}
    
    \s{
        Like the AC voltage source, the real AC source in Figure \ref{BildIdealStrom} has 
        an internal impedance $\underline{Z}_\mathrm{i}$, whereby this internal impedance is not connected in series as in the AC voltage source, 
        but rather in parallel to an ideal AC source. The resulting output current 
        $\underline{I}$ of the source depends on the connected load or the output voltage $\underline{U}$.
    }
    
    \b{
        AC source:
        \begin{itemize}
            \item<1-> Ideal AC source
            \item<2-> Real AC source with internal impedance $\underline{Z}_\mathrm{i}$
        \end{itemize}
    }
    
    \fu{
        \input{Tikz/src/Komplexe_Wechselstromrechnung/Quellen_Wechelgroessen2}
    }{{\bf Sources of alternating quantities.} Ideal AC source and ideal AC source. \label{BildIdealStrom}}
    
    \s{
        The alternating current of the ideal alternating current source is reduced by a portion due to the internal impedance. The 
        output current of the real alternating current source can be determined using equation \ref{GleichungQuel7}. Here, 
        starting from the source current $\underline{I}_\mathrm{0}$, the output voltage-dependent ($\underline{U}$) 
        alternating current component  $\underline{I}_\mathrm{i}$ is subtracted by the internal impedance $\underline{Z}_\mathrm{i}$.
    }
    
    \only<3>{
        \b{
            Output alternating current of a real alternating current source:
        }
        
        \begin{eq}
            \underline{I} = \underline{I}_\mathrm{0} - \underline{I}_\mathrm{i} = \underline{I}_\mathrm{0} - \frac{\underline{U}}{\underline{Z}_\mathrm{i}}     \label{GleichungQuel7}
        \end{eq}
    }
\end{frame}


\begin{frame}
    \renewcommand{\figurename}{Figure}
    \ftx{Equivalence of exchange sources}
    
    \s{
        {\bf AC conversion} \\
        
        
        Alternating current and alternating voltage sources can be converted into each other if both alternating sources have the same
        short-circuit current and the same open-circuit voltage.
    }
    
    \b{
        Conversion of alternating current sources and alternating voltage sources:
    }
    
    \fu{
        \input{Tikz/src/Komplexe_Wechselstromrechnung/Gegenueberstellung_reale_Wechselspannung_Wechselstromquelle}
    }{{\bf Comparison of a real alternating voltage source and a real alternating current source.} The real AC voltage source
        and the real AC current source are convertible into each other. \label{BildQuelleUmwandlung}}
    
    \s{
        Both AC sources are equivalent if, according to Equation \ref{GleichungQuel8} and Equation \ref{GleichungQuel9}, Ohm's law applies 
        with identical internal impedances for the AC voltage source and the AC current source. 
    }
    
    \b{
        Equivalent AC sources, with identical internal impedances:
    }
    
    \begin{eq}
        \underline{U}_\mathrm{0} = \underline{Z}_\mathrm{0} \cdot \underline{I}_\mathrm{0}      \label{GleichungQuel8}
    \end{eq}
    \begin{eq}
        \underline{Z}_\mathrm{i} = \underline{Z}_\mathrm{0}      \label{GleichungQuel9}
    \end{eq}
    
    \s{
        Here, the short-circuit current $\underline{I}_\mathrm{K}$ should be equal to the source current $\underline{I}_\mathrm{0}$ of the 
        alternating current source, which is identical to the output current $\underline{I}$. The output alternating voltage
        $\underline{U}$ of the alternating sources corresponds to the open-circuit voltage $\underline{U}_\mathrm{L}$.
    }
    
    \b{
        Source current of the current source equals the short-circuit current. \\
        Source voltage of the voltage source equals the open-circuit voltage.
    }
    
    \begin{eq}
        \underline{I} = \underline{I}_\mathrm{k} = \underline{I}_\mathrm{0} \qquad \text{und} \qquad \underline{U} = \underline{U}_\mathrm{L} = \underline{U}_\mathrm{0}     \label{GleichungQuel10}
    \end{eq}
    
\end{frame}



\begin{frame}
    \ftx{Example task}
    \begin{expl}{AC conversion}{Alternating current conversion:}
        
        \begin{enumerate}
            \only<1>{\item[]
                Given is a real AC voltage source with $\underline{U}_\mathrm{0} = 20\ kV \cdot \mathrm{e}^{\mathrm{j}30^o}$ and 
                $\underline{Z}_\mathrm{0} = (0.01 + j \cdot 2)\ \Omega$. Convert the alternating voltage source into an equivalent 
                alternating current source and calculate the short-circuit current.\\
                  }
                  \only<2>{
            \item[] 
                  \onslide<2>{
                      {\bf Short-circuit current:}
                      
                      
                      \begin{eq}
                          \underline{I}_\mathrm{0} = \frac{\underline{U}_\mathrm{0}}{\underline{Z}_\mathrm{0}}
                          = \frac{20\ kV \cdot \mathrm{e}^{\mathrm{j}30^o}}{(0,01 + j \cdot 2)\ \Omega}
                          = (5,04 + j \cdot 8,64)\ kA     \nonumber
                      \end{eq}
                  }
                  }
        \end{enumerate}
        
    \end{expl}
\end{frame}


\begin{frame}
    \renewcommand{\figurename}{Figure}
    \ftb{Complex voltage divider and complex current divider}
    
    \s{
        As with the consideration of direct current networks, the basic laws for analysing alternating current networks apply.
        The node rule and the mesh rule are to be applied equivalently here. The example node 
        and the example mesh already presented are shown once again in Figure \ref{BildKnotenMaschenKomplex} with complex alternating quantities.
        This results in the formulation of the following equation \ref{KnotenregelKomplex} as an 
        example of the complex node rule:
        
        \begin{eq}
            \sum_{k = 1}^{N}{\underline{I}_\mathrm{k}} = 0 = \underline{I}_\mathrm{1} + \underline{I}_\mathrm{2} - 
            \underline{I}_\mathrm{3} - \underline{I}_\mathrm{4} - \underline{I}_\mathrm{5}    \label{KnotenregelKomplex}
        \end{eq}
        
        As well as equation \ref{MaschengleichungKomplex} as an example of the complex mesh rule:
        
        \begin{eq}
            \sum_\mathrm{k = 1}^{N}{\underline{U}_\mathrm{k}} = \underline{U}_\mathrm{Z2} - \underline{U}_\mathrm{Z1}
            = 0  \label{MaschengleichungKomplex}
        \end{eq}
        
        In contrast to direct current analysis, where only real values are used in calculations, here it is only necessary to 
        ensure that complex alternating quantities are used. 
    }
    
    \b{
        As with the consideration of direct current networks, the basic laws apply to the analysis of alternating current networks.\\
        \begin{itemize}
            \item<2-> Knot rule:
                  \begin{eq}
                      \sum_{k = 1}^{N}{\underline{I}_\mathrm{k}} = 0 = \underline{I}_\mathrm{1} + \underline{I}_\mathrm{2} - 
                      \underline{I}_\mathrm{3} - \underline{I}_\mathrm{4} - \underline{I}_\mathrm{5}   \nonumber
                  \end{eq}
            \item<3-> Mesh rule:
                  \begin{eq}
                      \sum_\mathrm{k = 1}^{N}{\underline{U}_\mathrm{k}} = \underline{U}_\mathrm{Z2} - \underline{U}_\mathrm{Z1}
                      = 0  \nonumber
                  \end{eq}
        \end{itemize}
    }
    
    \fu{
        \resizebox{0.7\textwidth}{!}{
            \input{Tikz/src/Komplexe_Wechselstromrechnung/Knoten_und_Maschenregel}
        }
    }{{\bf Knot rule and mesh rule with complex alternating variables.} The node rule and the mesh rule behave
        identically with complex alternating quantities as with direct current networks. \label{BildKnotenMaschenKomplex}}
    
    
    
\end{frame}

\begin{frame}
    \renewcommand{\figurename}{Figure}
    \ftx{Complex voltage divider}
    
    \s{
        The complex voltage divider is used to reduce a total voltage, e.g. to adjust a signal level to the 
        input voltage range of a measuring circuit. Figure \ref{BildKomplxSpannungsteiler1} shows a complex 
        voltage divider consisting of two complex impedances. The complex total voltage $\underline{U}_\mathrm{0}$         
        is divided into the partial voltages $\underline{U}_\mathrm{1}$ and $\underline{U}_\mathrm{2}$ via the two 
        impedances $\underline{Z}_\mathrm{1}$ and $\underline{Z}_\mathrm{2}$.
        
        \fu{
            \resizebox{0.3\textwidth}{!}{
                \input{Tikz/src/Komplexe_Wechselstromrechnung/Komplexer_Spannungsteiler}
            }
        }{{\bf A complex voltage divider consisting of two impedances in series.} The complex total voltage
            $\underline{U}_\mathrm{0}$ is divided into the partial voltages $\underline{U}_\mathrm{1}$ and $\underline{U}_\mathrm{2}$
            \label{BildKomplxSpannungsteiler1}}
        
        
        The total voltage is divided in proportion to the complex impedances (see Equation \ref{GleichungTeiler1}). Based
        on the complex current $\underline{I}$ flowing through the entire series connection, the complex partial voltages
        $\underline{U}_\mathrm{1}$ and $\underline{U}_\mathrm{2}$ and their complex impedances $\underline{Z}_\mathrm{1}$ and 
        $\underline{Z}_\mathrm{2}$ can be determined in relation to the complex total voltage $\underline{U}_\mathrm{0}$ and the total complex
        impedance $\underline{Z}_\mathrm{1} + \underline{Z}_\mathrm{2}$ can be determined. 
        
        \begin{eq}
            \underline{I} = \frac{\underline{U}_\mathrm{0}}{\underline{Z}_\mathrm{1} + \underline{Z}_\mathrm{2}}
            = \frac{\underline{U}_\mathrm{1}}{\underline{Z}_\mathrm{1}}
            = \frac{\underline{U}_\mathrm{2}}{\underline{Z}_\mathrm{2}}     \label{GleichungTeiler1}
        \end{eq}
        
        By rearranging equation \ref{GleichungTeiler1}, the complex partial ratio $\underline{T}$ for equation 
        \ref{GleichungTeiler2} can be established: 
        
        \begin{eq}
            \underline{T} = \frac{\underline{U}_\mathrm{2}}{\underline{U}_\mathrm{0}}
            = \frac{\underline{Z}_\mathrm{2}}{\underline{Z}_\mathrm{1} + \underline{Z}_\mathrm{2}}    \label{GleichungTeiler2}
        \end{eq}
        
        Here, the ratio of the complex partial voltage $\underline{U}_\mathrm{2}$ and the complex total voltage 
        $\underline{U}_\mathrm{0}$ corresponds to the ratio of the complex impedance $\underline{Z}_\mathrm{2}$ and the total complex
        impedance of the series connection $\underline{Z}_\mathrm{1} + \underline{Z}_\mathrm{2}$. 
        The complex partial ratio can be established in the same way for the complex partial voltage $\underline{U}_\mathrm{1}$.
    }
    
    \b{
        \begin{minipage}[t]{0.4\textwidth}
            \fu{
                \resizebox{\textwidth}{!}{
                    \begin{circuitikz}
                        \draw (0,0) to [short, o-o] (3,0)
                        (0,4) to [short, o-] (1.5,4)
                        to [short,-o] (3,4)
                        (1.5,2) to [short, *-o] (3,2)
                        (1.5,4) to [R=$\underline{Z}_\mathrm{1}$, *-] (1.5,2)
                        to [R=$\underline{Z}_\mathrm{2}$, -*] (1.5,0);
                        \draw[-latex, thick, draw=voltage] (0,3.8) to node[left, color=voltage] {$\underline{U}_\mathrm{0}$} (0,0.2);
                        \draw[-latex, thick, draw=voltage] (3,3.8) to node[right, color=voltage] {$\underline{U}_\mathrm{1}$} (3,2.2);
                        \draw[-latex, thick, draw=voltage] (3,1.8) to node[right, color=voltage] {$\underline{U}_\mathrm{2}$} (3,0.2);
                        \pause
                        
                        \draw (0,4) to [short, i, name=I, o-] (1.5,4);
                        \iarrmore{I}{$\underline{I}$};     
                    \end{circuitikz}
                }
            }{}
        \end{minipage}
        \begin{minipage}[t]{0.5\textwidth}
            \begin{itemize}
                \item<1-> Reduction of a total voltage, e.g., to adjust a signal level to the
                        input voltage range of a measuring circuit.
                \item<2-> The total voltage is divided in proportion to the complex impedances:
                      \begin{eq}
                          \underline{I} = \frac{\underline{U}_\mathrm{0}}{\underline{Z}_\mathrm{1}+\underline{Z}_\mathrm{2}}
                          = \frac{\underline{U}_\mathrm{1}}{\underline{Z}_\mathrm{1}}
                          = \frac{\underline{U}_\mathrm{2}}{\underline{Z}_\mathrm{2}}     \nonumber
                      \end{eq}
                \item<3-> Complex divisor ratio
                      \begin{eq}
                          \underline{T} = \frac{\underline{U}_\mathrm{2}}{\underline{U}_\mathrm{0}}
                          = \frac{\underline{Z}_\mathrm{2}}{\underline{Z}_\mathrm{1}+\underline{Z}_\mathrm{2}}    \nonumber
                      \end{eq}
            \end{itemize}
        \end{minipage}
    }
\end{frame}


\begin{frame}
    \renewcommand{\figurename}{Figure}
    \ftx{Complex voltage divider: Application example high-voltage technology}
    
    \s{
        The complex voltage divider is used, for example, in ohmic-capacitive dividers in high-voltage technology. 
        Figure \ref{BildKomplxSpannungsteiler2} shows four different dividers that are used, for example, in 
        high-voltage technology. 
        
        \fu{
            \resizebox{0.5\textwidth}{!}{
                \input{Tikz/src/Komplexe_Wechselstromrechnung/Gegenueberstellung_Teiler_Hochspannung}
            }
        }{{\bf Comparison of different dividers in high-voltage technology.} Ohmic divider for use with direct voltage and
        capacitive divider and damped capacitive divider for use with alternating voltage. The ohmic-capacitive divider can be used for 
        both types of voltage. \label{BildKomplxSpannungsteiler2}}
        
        In addition to the ohmic-capacitive divider, the ohmic divider, the 
        capacitive divider and the damped capacitive divider are also presented. The ohmic divider consists purely of ohmic resistors, 
        which are connected in series, and the capacitive divider consists purely of capacitors. 
        The damped capacitive divider consists of several modules with an upstream coil. The modules are composed of series connections of an 
        resistive resistor and a capacitor.  
    }
    
    \b{
        \begin{minipage}[t]{0.55\textwidth}
            \fu{
                \resizebox{\textwidth}{!}{
                    \begin{circuitikz}
                        \draw (0,10) to [short, o-] (1,10)
                        to [R] (1,8) to [R] (1,6) to [R] (1,4) to [R] (1,2) to [R] (1,0)
                        (1,2) to [short, *-o] (2,2)
                        (0,0) to [short, o-*] (1,0) to [short, -o] (2,0);
                        \draw[-latex, thick, draw=voltage] (0,9.8) to node[left, color=voltage] {$\underline{U}_\mathrm{E}$} (0,0.2);
                        \draw[-latex, thick, draw=voltage] (2,1.8) to node[right, color=voltage] {$\underline{U}_\mathrm{A}$} (2,0.2);
                        \draw (1,0) node [label=below:Ohmscher Teiler]{};
                        \pause
                        
                        \draw (4,10) to [short, o-*] (5,10)
                        to [R] (5,8) to [R] (5,6) to [R] (5,4) to [R] (5,2) to [R] (5,0)
                        (5,10) to [short] (6,10)
                        to [C] (6,8) to [C] (6,6) to [C] (6,4) to [C] (6,2) to [C] (6,0)
                        (5,8) to [short, *-*] (6,8)
                        (5,6) to [short, *-*] (6,6)
                        (5,4) to [short, *-*] (6,4)
                        (5,2) to [short, *-*] (6,2) to [short, -o] (7,2)
                        (4,0) to [short, o-o] (7,0) (5,0) to [short, *-*] (6,0);
                        \draw[-latex, thick, draw=voltage] (4,9.8) to node[left, color=voltage] {$\underline{U}_\mathrm{E}$} (4,0.2);
                        \draw[-latex, thick, draw=voltage] (7,1.8) to node[right, color=voltage] {$\underline{U}_\mathrm{A}$} (7,0.2);
                        \draw (5.5,0) node [label=below:Ohmsch-kapazitiver Teiler]{};
                        \pause
                        
                        \draw (9,10) to [short, o-] (10,10)
                        to [C] (10,8) to [C] (10,6) to [C] (10,4) to [C] (10,2) to [C] (10,0)
                        (10,2) to [short, *-o] (11,2)
                        (9,0) to [short, o-*] (10,0) to [short, -o] (11,0);
                        \draw[-latex, thick, draw=voltage] (9,9.8) to node[left, color=voltage] {$\underline{U}_\mathrm{E}$} (9,0.2);
                        \draw[-latex, thick, draw=voltage] (11,1.8) to node[right, color=voltage] {$\underline{U}_\mathrm{A}$} (11,0.2);
                        \draw (10,0) node [label=below:Kapazitiver Teiler]{};
                        \pause
                        
                        \draw (13,10) to [short, o-] (13.5,10) to [L] (14.5,10) to [short] (15,10)
                        to [R] (15,8.6) to [C] (15,8) to [R] (15,6.6) to [C] (15,6) to [R] (15,4.6) to [C] (15,4) to [R] (15,2.6) to [C] (15,2) to [R] (15,0.6) to [C] (15,0)
                        (15,2) to [short, *-o] (16,2) 
                        (13,0) to [short, o-*] (15,0) to [short, -o] (16,0);
                        \draw[-latex, thick, draw=voltage] (13,9.8) to node[left, color=voltage] {$\underline{U}_\mathrm{E}$} (13,0.2);
                        \draw[-latex, thick, draw=voltage] (16,1.8) to node[right, color=voltage] {$\underline{U}_\mathrm{A}$} (16,0.2);
                        \draw (14.5,0) node [label=below:Gedämpft kapazitiver Teiler]{};
                    \end{circuitikz}
                }
            }{}
        \end{minipage}
        \begin{minipage}[t]{0.35\textwidth}
            \begin{itemize}
                \item<1-> Ohmic divider (DC voltage)
                \item<2-> Ohmic-capacitive divider (DC and AC voltage)
                \item<3-> Capacitive divider (AC voltage)
                \item<4-> Damped capacitive divider (AC voltage)
            \end{itemize}
        \end{minipage}
    }
\end{frame}

\begin{frame}
    \renewcommand{\figurename}{Figure}
    \ftx{Complex voltage divider}
    
    \s{
        A simple example of such an ohmic-capacitive divider
        is shown in Figure \ref{BildKomplxSpannungsteiler3}. The ohmic-capacitive divider consists of at least two 
        parallel circuits connected in series, each comprising an ohmic resistor and a capacitor, known as an RC element .
        In Figure \ref{BildKomplxSpannungsteiler3}, the first RC element consists of the resistor $R_1$ and the capacitor $C_1$. The second RC element consists of the resistor $R_2$ and the capacitor $C_2$. The complex total voltage         
        $\underline{U}_\mathrm{0}$ is applied across both RC elements. The complex partial voltage $\underline{U}_\mathrm{2}$ is applied across 
        the second RC element. The partial ratio therefore has ohmic and capacitive components. The capacitance of the 
        capacitor results in a frequency-dependent partial ratio.
        
        \fu{
            \resizebox{0.3\textwidth}{!}{
                \input{Tikz/src/Komplexe_Wechselstromrechnung/Spannungsteiler_ohmsch-kapazitiv}
            }
        }{{\bf An ohmic-capacitive voltage divider consisting of two RC elements.} The RC elements each consist of a parallel connection
        of a resistor and a capacitor. \label{BildKomplxSpannungsteiler3}}
        
        A special condition exists when the complex voltage divider is balanced. This condition describes the identical 
        ratio of the components to each other. This condition is described by equation \ref{GleichungTeiler3}. Here 
        the product of the resistor $R_1$ and the capacitor $C_1$ must be identical to the product of the resistor $R_2$ and the 
        capacitor $C_2$. In practice, this balancing condition can be achieved using variable components such as trimmer capacitors. 
        
        \begin{eq}
            R_1 \cdot C_1 = R_2 \cdot C_2   \label{GleichungTeiler3}
        \end{eq}
        
    }
    
    \b{
        \begin{minipage}[t]{0.4\textwidth}
            \fu{
                \resizebox{\textwidth}{!}{
                    \begin{circuitikz}
                        \draw (0,0) to [short, o-*] (2.25,0)
                        to [short, -o] (4.5,0)
                        (0,6) to [short, o-] (2.25,6) to [short] (2.25,5.5)
                        (1.25,5.5) to [short] (2.25,5.5) to [short, *-] (3.25,5.5)
                        (1.25,3.5) to [short] (2.25,3.5) to [short, *-] (3.25,3.5)
                        (1.25,2.5) to [short] (2.25,2.5) to [short, *-] (3.25,2.5)
                        (1.25,0.5) to [short] (2.25,0.5) to [short, *-] (3.25,0.5)
                        (2.25,3.5) to [short] (2.25,2.5)
                        (2.25,0) to [short] (2.25,0.5)
                        (2.25,3) to [short, *-o] (4.5,3)
                        (1.25,5.5) to [R=$R_\mathrm{1}$] (1.25,3.5)
                        (3.25,5.5) to [C=$C_\mathrm{1}$] (3.25,3.5)
                        (1.25,2.5) to [R=$R_\mathrm{2}$] (1.25,0.5)
                        (3.25,2.5) to [C=$C_\mathrm{2}$] (3.25,0.5);
                        \draw[-latex, thick, draw=voltage] (0,5.8) to node[left, color=voltage] {$\underline{U}_\mathrm{0}$} (0,0.2);
                        \draw[-latex, thick, draw=voltage] (4.5,2.8) to node[right, color=voltage] {$\underline{U}_\mathrm{2}$} (4.5,0.2);
                    \end{circuitikz}
                }
            }{}
        \end{minipage}
        \begin{minipage}[t]{0.5\textwidth}
            \begin{itemize}
                \item<1-> Ohmic and capacitive division ratio
                \item<2-> Impedances dependent on frequency
                \item<3-> Complex division ratio is frequency-dependent
                \item<4-> The division ratio is only in exceptional cases
                independent of frequency!
                      \begin{eq}
                          R_1 \cdot C_1 = R_2 \cdot C_2   \nonumber    
                      \end{eq}
            \end{itemize}
        \end{minipage}
    }
\end{frame}

\begin{frame}
    \renewcommand{\figurename}{Figure}
    \ftx{Complex voltage divider: Application example probe head}
    
    \s{
        Another example of an ohmic-capacitive divider is the probe of an oscilloscope. Such a
        divider is shown in Figure \ref{BildKomplxSpannungsteiler4}. 
        In principle, the image does not differ significantly from the ohmic-capacitive divider. However, 
        here the first and second RC elements are spatially separated from each other 
        and enclosed in a housing. The first RC element represents the probe and the second 
        RC element represents the circuitry in the oscilloscope. With the ohmic division ratio:

        \begin{eq}
            \frac{\underline{U}_\mathrm{A}}{\underline{U}_\mathrm{E}} = \frac{R_\mathrm{2}}{R_\mathrm{1}+R_\mathrm{2}}
        \end{eq}
        
        and the capacitive ratio:
        \begin{eq}
            \frac{\underline{U}_\mathrm{A}}{\underline{U}_\mathrm{E}} = \frac{C_\mathrm{2}}{C_\mathrm{1}+C_\mathrm{2}}
        \end{eq}
        
        To calibrate the probe, the ohmic and capacitive division ratios can be equated, as the 
        division ratios must be identical:
        \begin{eq}
            \frac{R_\mathrm{2}}{R_\mathrm{1}+R_\mathrm{2}} = \frac{C_\mathrm{2}}{C_\mathrm{1}+C_\mathrm{2}}
        \end{eq}
        
        The adjustment condition according to equation \ref{GleichungTeiler3} still applies. 
        The divider is adjusted via the adjustable capacitance $C_1$.
 
        
        \fu{
            \resizebox{0.3\textwidth}{!}{
               \input{Tikz/src/Komplexe_Wechselstromrechnung/Spannungsteiler_ohmsch-kapazitiv2}
            }
        }{{\bf An ohmic-capacitive voltage divider as the probe head of an oscilloscope.} The complex voltage divider is balanced via capacitor $C_1$.
        The voltage divider is balanced via capacitor $C_1$. \label{BildKomplxSpannungsteiler4}}
        
    }
    
    \b{
        \begin{minipage}[t]{0.5\textwidth}
            \fu{
                \resizebox{\textwidth}{!}{
                    \begin{circuitikz}
                        \draw (0,0) to [short, o-o] (4.5,0)
                        (0,6) to [short, o-] (2.25,6) to [short] (2.25,5.5)
                        (1.25,5.5) to [short] (2.25,5.5) to [short, *-] (3.25,5.5)
                        (1.25,3.5) to [short] (2.25,3.5) to [short, *-] (3.25,3.5)
                        (1.25,2.5) to [short] (2.25,2.5) to [short, *-] (3.25,2.5)
                        (1.25,0.5) to [short] (2.25,0.5) to [short, *-] (3.25,0.5)
                        (2.25,3.5) to [short] (2.25,2.5)
                        (2.25,0) to [short, *-] (2.25,0.5)
                        (2.25,3) to [short, *-o] (4.5,3)
                        (1.25,5.5) to [R=$R_\mathrm{1}$] (1.25,3.5)
                        (3.25,5.5) to [C=$C_\mathrm{1}$] (3.25,3.5)
                        (1.25,2.5) to [R=$R_\mathrm{2}$] (1.25,0.5)
                        (3.25,2.5) to [C=$C_\mathrm{2}$] (3.25,0.5);
                        \draw[-latex, thick, draw=voltage] (0,5.8) to node[left, color=voltage] {$\underline{U}_\mathrm{E}$} (0,0.2);
                        \draw[-latex, thick, draw=voltage] (4.5,2.8) to node[right, color=voltage] {$\underline{U}_\mathrm{A}$} (4.5,0.2);
                        
                        \draw [dashed] (0.5,5.75) rectangle (4.25,3.25);
                        \draw (3.25,5.5) node [label=above:Probe]{};
                        
                        \draw [dashed] (0.5,2.75) rectangle (4.25,0.25);
                        \draw (1.25,2.5) node [label=above:Oscilloscope]{};
                    \end{circuitikz}
                }
            }{}
        \end{minipage}
        \begin{minipage}[t]{0.45\textwidth}
            \begin{itemize}
                \item<1-> Ohmic ratio:
                      \begin{eq}
                          \frac{\underline{U}_\mathrm{A}}{\underline{U}_\mathrm{E}} = \frac{R_\mathrm{2}}{R_\mathrm{1}+R_\mathrm{2}}   \nonumber    
                      \end{eq}
                \item<2-> Capacitive partial ratio:
                      \begin{eq}
                          \frac{\underline{U}_\mathrm{A}}{\underline{U}_\mathrm{E}} = \frac{C_\mathrm{2}}{C_\mathrm{1}+C_\mathrm{2}}   \nonumber  
                      \end{eq}
                \item<3-> Equivalence of ohmic and capacitive partial ratios:
                      \begin{eq}
                          \frac{R_\mathrm{2}}{R_\mathrm{1}+R_\mathrm{2}} = \frac{C_\mathrm{2}}{C_\mathrm{1}+C_\mathrm{2}}   \nonumber    
                      \end{eq}
                \item<4-> Comparison condition:
                      \begin{eq}
                          R_2 \cdot C_2 = R_1 \cdot C_1   \nonumber    
                      \end{eq}
            \end{itemize}
        \end{minipage}
    }
\end{frame}


\begin{frame}
    \renewcommand{\figurename}{Figure}
    \ftx{Complex current divider}
    
    \s{
        The complex current divider, as shown in Figure \ref{BildKomplxStromteiler}, divides the complex 
        total current $\underline{I}$ into the partial currents $\underline{I}_\mathrm{1}$ and $\underline{I}_\mathrm{2}$. The two
        partial currents $\underline{I}_\mathrm{1}$ and $\underline{I}_\mathrm{2}$ flow through the admittances $\underline{Y}_\mathrm{1}$
        and $\underline{Y}_\mathrm{2}$.
        
        \fu{
            \resizebox{0.3\textwidth}{!}{
                \input{Tikz/src/Komplexe_Wechselstromrechnung/Stromteiler_komplex}
            }
        }{{\bf A complex current divider consisting of two admittances.} The complex current $\underline{I}$ is divided
        into the currents $\underline{I}_\mathrm{1}$ and $\underline{I}_\mathrm{2}$. \label{BildKomplxStromteiler}}
        
        The complex voltage $\underline{U}$ can be calculated according to Equation \ref{GleichungStromteiler1} using the ratio of 
        the total current $\underline{I}$ and the total admittance $\underline{Y}_\mathrm{1} + \underline{Y}_\mathrm{2}$ or 
        using the ratios of the individual current branches:
        
        \begin{eq}
            \underline{U} = \frac{\underline{I}_\mathrm{0}}{\underline{Y}_\mathrm{1}+\underline{Y}_\mathrm{2}} = 
            \frac{\underline{I}_\mathrm{1}}{\underline{Y}_\mathrm{1}} = 
            \frac{\underline{I}_\mathrm{2}}{\underline{Y}_\mathrm{2}}   \label{GleichungStromteiler1}
        \end{eq}
        
        The division ratio can be defined for the two current branches according to equation \ref{GleichungStromteiler2}: 

        \begin{eqa}
            \underline{T}_\mathrm{i1} = \frac{\underline{I}_\mathrm{1}}{\underline{I}_\mathrm{0}} = 
            \frac{\underline{Y}_\mathrm{1}}{\underline{Y}_\mathrm{1}+\underline{Y}_\mathrm{2}}  \label{GleichungStromteiler2}  \\ 
            \underline{T}_\mathrm{i2} = \frac{\underline{I}_\mathrm{2}}{\underline{I}_\mathrm{0}} = 
            \frac{\underline{Y}_\mathrm{2}}{\underline{Y}_\mathrm{1}+\underline{Y}_\mathrm{2}} \label{GleichungStromteiler3}
        \end{eqa}
        
    }
    
    \b{
        \begin{minipage}[t]{0.4\textwidth}
            \fu{
                \resizebox{\textwidth}{!}{
                    \begin{circuitikz}
                        \draw (0,0) to [short, o-] (2.25,0)
                        to [short, -*] (2.25,0.5)
                        (1.25,0.5) to [short] (3.25,0.5) 
                        (1.25,0.5) to [R=$\underline{Y}_\mathrm{1}$, i_<, name=I1] (1.25,3.5)
                        (3.25,0.5) to [R=$\underline{Y}_\mathrm{2}$, i_<, name=I2] (3.25,3.5)
                        (1.25,3.5) to [short] (3.25,3.5)
                        (2.25,3.5) to [short, *-] (2.25,4)
                        (0,4) to [short, i, name=I, o-] (2.25,4);
                        \draw[-latex, thick, draw=voltage] (4,3.3) to node[right, color=voltage] {$\underline{U}$} (4,0.7);
                        
                        \iarrmore{I}{$\underline{I}_\mathrm{0}$}; 
                        \iarrmore{I1}{$\underline{I}_\mathrm{1}$};
                        \iarrmore{I2}{$\underline{I}_\mathrm{2}$};  
                    \end{circuitikz}
                }
            }{}
        \end{minipage}
        \begin{minipage}[t]{0.5\textwidth}
            \begin{itemize}
                \item<1-> The total current is divided in proportion to the complex admittances:
                      \begin{eq}
                          \underline{U} = \frac{\underline{I}_\mathrm{0}}{\underline{Y}_\mathrm{1}+\underline{Y}_\mathrm{2}} = 
                          \frac{\underline{I}_\mathrm{1}}{\underline{Y}_\mathrm{1}} = 
                          \frac{\underline{I}_\mathrm{2}}{\underline{Y}_\mathrm{2}}   \nonumber    
                      \end{eq}
                \item<2-> Impedances depending on frequency
                      \begin{eqa}
                          \underline{T}_\mathrm{i1} = \frac{\underline{I}_\mathrm{1}}{\underline{I}_\mathrm{0}} = 
                          \frac{\underline{Y}_\mathrm{1}}{\underline{Y}_\mathrm{1}+\underline{Y}_\mathrm{2}}  \\ 
                          \underline{T}_\mathrm{i2} = \frac{\underline{I}_\mathrm{2}}{\underline{I}_\mathrm{0}} = 
                          \frac{\underline{Y}_\mathrm{2}}{\underline{Y}_\mathrm{1}+\underline{Y}_\mathrm{2}} \nonumber    
                      \end{eqa}
            \end{itemize}
        \end{minipage}
    }
    
\end{frame}

\newpage